% META: {"id": "1NSI-TABLES-RE-C2", "chapitre": "1NSI-TABLES", "type_objet": "remediation", "capacites": ["P-TABLE-02"], "status": "generated"}

%---------------------------------------------------------------
% FICHE DE REMEDIATION — C2 : doublons sur une colonne, pas sur la ligne entiere
%---------------------------------------------------------------

\begin{exercice}{1NSI-TAB-RE-C2-EX1}{1}{10}
Un élève cherche les élèves inscrits deux fois (même nom) dans la table suivante, mais
compare des \textbf{lignes entières} au lieu de la seule colonne \lstinline{nom} :
\begin{python}
eleves = [
    {"nom": "Ali", "age": 16, "classe": "1A"},
    {"nom": "Ali", "age": 17, "classe": "1B"},
]
lignes_str = [str(e) for e in eleves]
doublons_naifs = [chaine for chaine in set(lignes_str) if lignes_str.count(chaine) > 1]
\end{python}
\begin{enumerate}
  \item Que vaut \lstinline{doublons_naifs} ? Le nom \lstinline{"Ali"} est-il pourtant présent deux fois dans la table ?
  \item Pourquoi la méthode de l'élève échoue-t-elle à détecter ce cas ?
  \item Réécrire une détection de doublons qui compare uniquement la colonne \lstinline{nom}, et vérifier qu'elle détecte bien \lstinline{"Ali"}.
\end{enumerate}
\end{exercice}

% BEGIN-VERIFY
% eleves = [
%     {"nom": "Ali", "age": 16, "classe": "1A"},
%     {"nom": "Ali", "age": 17, "classe": "1B"},
% ]
% lignes_str = [str(ligne) for ligne in eleves]
% doublons_naifs = [chaine for chaine in set(lignes_str) if lignes_str.count(chaine) > 1]
% assert doublons_naifs == []
% noms = [ligne["nom"] for ligne in eleves]
% doublons_colonne = [n for n in set(noms) if noms.count(n) > 1]
% assert doublons_colonne == ["Ali"]
% END-VERIFY

\begin{corrige}{1NSI-TAB-RE-C2-EX1}
\textbf{1.} \lstinline{doublons_naifs} vaut \lstinline{[]} (liste vide) : aucun doublon
n'est détecté. Pourtant, \lstinline{"Ali"} apparaît bien deux fois dans la colonne
\lstinline{nom}.

\textbf{2.} Les deux lignes ont le même \lstinline{nom} mais des valeurs
\textbf{différentes} pour \lstinline{age} et \lstinline{classe} : converties en chaînes
avec \lstinline{str(...)}, ce sont donc deux chaînes \textbf{différentes}, qui n'apparaissent
chacune qu'une seule fois. La comparaison ligne entière ne peut détecter que des lignes
strictement identiques sur \textbf{toutes} les colonnes, pas un doublon sur une seule
colonne.

\textbf{3.}
\begin{python}
noms = [ligne["nom"] for ligne in eleves]
doublons_colonne = [n for n in set(noms) if noms.count(n) > 1]
\end{python}
\lstinline{doublons_colonne} vaut \lstinline{["Ali"]} : cette fois, le doublon est bien
détecté, car on compare uniquement la colonne concernée.
\end{corrige}
